\documentclass[11pt]{article}

% ---------- packages ----------
\usepackage[utf8]{inputenc}
\usepackage[T1]{fontenc}
\usepackage[a4paper,margin=2.5cm]{geometry}
\usepackage{graphicx}
\usepackage{amsmath,amssymb}
\usepackage{booktabs}
\usepackage{array}
\usepackage{float}
\usepackage{caption}
\usepackage{siunitx}
\usepackage[hidelinks]{hyperref}

\captionsetup{font=small,labelfont=bf}
\renewcommand{\arraystretch}{1.15}

% ---------- title page ----------
\begin{document}
\begin{titlepage}
\centering
% If you add the logo file (sapienza-logo.png) it will be included automatically.
\IfFileExists{sapienza-logo.png}{%
  \includegraphics[width=0.32\textwidth]{sapienza-logo.png}\\[0.4cm]
}{\vspace*{1.5cm}}

{\scshape\large Dipartimento di Ingegneria Informatica, Automatica\\ e Gestionale\par}
\vspace{2.5cm}

{\huge\bfseries Parameter-Efficient Few-Step Distillation\\[0.2cm] of Diffusion Models via LoRA\par}
\vspace{0.4cm}
{\scshape\Large Computer Vision --- Project 5\par}
\vspace{0.2cm}
{\itshape\large Lightweight Diffusion Models for Resource-Constrained Environments\par}
\vspace{2.5cm}

\begin{minipage}{0.45\textwidth}
\begin{flushleft}\large
\textbf{Professor:}\\
Irene Amerini
\end{flushleft}
\end{minipage}
\begin{minipage}{0.45\textwidth}
\begin{flushright}\large
\textbf{Students:}\\
Claudia Giuseppini --- 2069387\\
Benedetta Conti --- 2045621\\
Fabio Falconi --- 2048079
\end{flushright}
\end{minipage}

\vfill
\rule{\textwidth}{0.4pt}\\[0.2cm]
{\large Academic Year 2025/2026\par}
\end{titlepage}

\tableofcontents
\newpage

% ==================================================================
\section{Introduction}
% ==================================================================
Denoising Diffusion Probabilistic Models (DDPMs)~\cite{ho2020ddpm} produce
high-quality images, but their sampling procedure is an iterative Markov chain of
hundreds to thousands of network evaluations, which makes inference slow and
expensive. This is a critical limitation for deployment on edge devices or for
training in resource-constrained environments.

This project, following the Project~5 specification, implements a standard DDPM
baseline on CIFAR-10 and integrates a methodological novelty aimed at accelerating
sampling. We then rigorously evaluate the trade-off between computational efficiency
(sampling time) and generation quality (FID, Inception Score) as well as sample
diversity (Precision/Recall).

Deterministic samplers such as DDIM~\cite{song2021ddim} reduce the number of steps,
but quality collapses once the step count becomes very small. The root cause is the
\emph{accumulated discretization error} along the reverse trajectory, not a failure
of noise prediction at individual timesteps. Our contribution directly targets this
error while keeping the adaptation cheap:

\begin{quote}
\textbf{Parameter-efficient few-step distillation via LoRA.} We freeze the trained
DDPM (the \emph{teacher}) and train only low-rank adapters~\cite{hu2021lora} injected
into the attention layers (the \emph{student}) with a progressive-distillation-style
trajectory-matching objective~\cite{salimans2022progressive}. The student learns to
reproduce, in very few steps, what the teacher produces with many DDIM sub-steps.
\end{quote}

Because only the adapters are trained (a small fraction of the weights) and they can
be merged into the base weights at inference, the method adds essentially no inference
overhead while providing large speed-ups over the full-step baseline.

% ==================================================================
\section{Related Work and Identified Limitations}
% ==================================================================
\begin{itemize}
  \item \textbf{DDPM}~\cite{ho2020ddpm} --- our baseline generative model.
        \emph{Limitation:} very slow sampling ($T=1000$ sequential evaluations).
  \item \textbf{DDIM}~\cite{song2021ddim} --- deterministic non-Markovian sampler
        enabling fewer steps. \emph{Limitation:} quality collapses under aggressive
        step reduction (in our measurements, FID rises from $\sim\!29$ at 100 steps to
        $>\!100$ at 4 steps).
  \item \textbf{Progressive Distillation}~\cite{salimans2022progressive} --- distils a
        teacher into a student that halves the number of steps. \emph{Limitation:}
        requires full-model fine-tuning at each stage, costly on constrained hardware.
  \item \textbf{LoRA}~\cite{hu2021lora} --- low-rank adapters for parameter-efficient
        adaptation.
\end{itemize}

Our novelty sits at the intersection: we bring progressive-distillation-style
trajectory matching to diffusion \emph{but perform it with LoRA only}, addressing the
cost limitation of distillation while attacking the discretization-error limitation of
few-step DDIM.

% ==================================================================
\section{Method}
% ==================================================================
\subsection{Baseline DDPM}
We use an attention U-Net (encoder--decoder with skip connections, GroupNorm,
self-attention at low resolution, sinusoidal timestep embeddings). The forward
diffusion process noises a clean image $x_0$ as
\begin{equation}
  x_t = \sqrt{\bar{\alpha}_t}\,x_0 + \sqrt{1-\bar{\alpha}_t}\,\varepsilon,
  \qquad \varepsilon \sim \mathcal{N}(0, I),
  \label{eq:forward}
\end{equation}
where $\bar{\alpha}_t = \prod_{s=1}^{t}(1-\beta_s)$ follows a linear $\beta$-schedule
with $T=1000$. The network $\varepsilon_\theta$ is trained with the standard
noise-prediction objective
\begin{equation}
  \mathcal{L}_{\mathrm{DDPM}} =
  \mathbb{E}_{x_0,\,t,\,\varepsilon}
  \big\| \varepsilon - \varepsilon_\theta(x_t, t) \big\|_2^2 ,
  \label{eq:ddpm-loss}
\end{equation}
using an exponential moving average (EMA) of the weights and mixed precision (AMP).

\subsection{Fast Sampling (DDIM)}
Deterministic DDIM ($\eta=0$) advances from timestep $t$ to $t'$ via
\begin{equation}
  \hat{x}_0 = \frac{x_t - \sqrt{1-\bar{\alpha}_t}\,\varepsilon_\theta(x_t,t)}
                   {\sqrt{\bar{\alpha}_t}},
  \qquad
  x_{t'} = \sqrt{\bar{\alpha}_{t'}}\,\hat{x}_0
           + \sqrt{1-\bar{\alpha}_{t'}}\,\varepsilon_\theta(x_t,t).
  \label{eq:ddim}
\end{equation}
Reducing the number of visited timesteps accelerates sampling but increases the
discretization error between consecutive $x_{t'}$ estimates.

\subsection{Novelty: LoRA Few-Step Distillation}
The trained baseline (EMA weights) is the frozen \emph{teacher} $\varepsilon_\phi$.
The \emph{student} $\varepsilon_\theta$ is the same network with LoRA adapters injected
into the attention projections; only the adapters are trainable.

For a target step budget $K$ we build a short schedule and, for a randomly sampled
interval $(t \rightarrow t')$, we forward-diffuse a real image to $x_t$ via
Eq.~\eqref{eq:forward}. The teacher refines $x_t$ from $t$ to $t'$ using
$m$ small DDIM sub-steps (Eq.~\eqref{eq:ddim}), producing the target
$x_{t'}^{\text{teacher}}$; the student must reach the same point in a
\emph{single} DDIM step, producing $x_{t'}^{\text{student}}$. The objective is a
trajectory-matching loss with an auxiliary noise-consistency term:
\begin{equation}
  \mathcal{L}_{\mathrm{distill}} =
  \big\| x_{t'}^{\text{student}} - x_{t'}^{\text{teacher}} \big\|_2^2
  \;+\; \lambda\,
  \big\| \varepsilon_\theta(x_t,t) - \varepsilon_\phi(x_t,t) \big\|_2^2 .
  \label{eq:distill}
\end{equation}
With $m=2$ this reduces to canonical progressive distillation. We train one student per
target $K \in \{8,4,2,1\}$. At inference the adapters can be merged into the base
weights, so the distilled model costs the same as the base model at the same $K$.

% ==================================================================
\section{Experimental Setup}
% ==================================================================
\begin{itemize}
  \item \textbf{Data:} CIFAR-10 ($32\times32$ RGB); metrics computed against the
        CIFAR-10 test split via \texttt{torch-fidelity}.
  \item \textbf{Model:} U-Net with \texttt{base\_channels}$=64$,
        \texttt{time\_dim}$=256$ (single-GPU scale).
  \item \textbf{Training:} 80 epochs, batch size 128, AdamW, EMA, AMP.
  \item \textbf{Distillation:} 10 epochs per student, LoRA on the attention projections
        (rank 4 by default; ranks 2 and 8 for the ablation).
  \item \textbf{Metrics:} FID, KID, Inception Score (fidelity); Precision/Recall
        (diversity)~\cite{kynkaanniemi2019precision}; seconds per image (speed).
        All computed on \num{10000} generated images from a shared, seeded noise bank.
  \item \textbf{Reproducibility:} fixed seed (42), deterministic algorithms, seeded
        \texttt{torch.Generator}, pinned package versions.
\end{itemize}

% ==================================================================
\section{Results}
% ==================================================================
\subsection{Baseline: quality and speed vs.\ number of steps}
Table~\ref{tab:steps} reports the baseline across step budgets. Fewer steps means
faster but lower-quality sampling; crucially, \textbf{recall (diversity) collapses}
long before precision does.

\begin{table}[H]
\centering
\caption{Baseline: DDIM at different step counts (FID on \num{10000} images).
The DDPM row uses \num{1000} images and is reported only for speed-up factors.}
\label{tab:steps}
\begin{tabular}{llrrrrrr}
\toprule
Sampler & Steps & FID $\downarrow$ & IS $\uparrow$ & Precision & Recall & s/img & Speed-up$^\dagger$ \\
\midrule
DDPM (ref) & 1000 & 67.3$^{*}$ & 6.26 & 0.847 & 0.381 & 0.678 & $1\times$ \\
DDIM & 100 & 29.33 & 7.23 & 0.751 & 0.326 & 0.053 & $12.8\times$ \\
DDIM & 50  & 29.79 & 7.32 & 0.777 & 0.306 & 0.027 & $25.2\times$ \\
DDIM & 20  & 32.43 & 7.32 & 0.834 & 0.243 & 0.011 & $60.5\times$ \\
DDIM & 10  & 40.89 & 7.02 & 0.893 & 0.141 & 0.006 & $112.9\times$ \\
DDIM & 8   & 47.12 & 6.83 & 0.903 & 0.102 & 0.005 & $136.5\times$ \\
DDIM & 4   & 103.28 & 5.29 & 0.853 & 0.015 & 0.003 & $242.4\times$ \\
DDIM & 2   & 188.84 & 3.73 & 0.401 & 0.001 & 0.002 & $\sim\!360\times$ \\
DDIM & 1   & 325.11 & 1.56 & 0.005 & 0.000 & 0.001 & $\sim\!510\times$ \\
\bottomrule
\end{tabular}

\vspace{2pt}
{\footnotesize $^\dagger$ vs.\ DDPM-1000. \quad $^{*}$ On \num{1000} images (noisier).}
\end{table}

\paragraph{Fidelity vs.\ Diversity.} As steps drop from 100 to 8, precision actually
\emph{rises} ($0.75\!\rightarrow\!0.90$) while recall \emph{collapses}
($0.33\!\rightarrow\!0.10$, and to $\sim\!0$ at 4 steps and below). Few-step sampling
does not primarily degrade individual samples --- it shrinks the portion of the data
distribution the model covers. FID alone would hide this; Precision/Recall makes it
explicit.

\subsection{Novelty: matched-step comparison}
Table~\ref{tab:paired} compares plain DDIM-$K$ against the distilled LoRA-$K$ at the
same step budget, same initial noise, and same base weights --- so the only difference
is the LoRA adapter.

\begin{table}[H]
\centering
\caption{Matched step budget $K$: plain DDIM vs.\ distilled LoRA.}
\label{tab:paired}
\begin{tabular}{crrrrrc}
\toprule
$K$ & DDIM FID & Distilled FID & $\Delta$FID & DDIM Recall & Distilled Recall & Time ratio \\
\midrule
8 & 47.12  & \textbf{38.24} & \textbf{+8.9}  & 0.102 & 0.159 & $1.01\times$ \\
4 & 103.28 & \textbf{68.34} & \textbf{+34.9} & 0.015 & 0.051 & $1.00\times$ \\
2 & 188.84 & \textbf{159.12}& \textbf{+29.7} & 0.001 & 0.002 & $1.01\times$ \\
1 & 325.11 & 326.57 & $-1.5$ & 0.000 & 0.000 & $1.00\times$ \\
\bottomrule
\end{tabular}
\end{table}

The distilled student improves FID at $K=8,4,2$, with the largest gain at
\textbf{$K=4$ ($+34.9$)}, and partially \textbf{restores diversity} (recall $+0.057$ at
$K=8$). At $K=1$ both models fail. Inference cost is unchanged (time ratio $\approx 1$):
the distilled $K=8$ model is $\sim\!10.6\times$ faster than DDIM-100 and the $K=4$ model
$\sim\!18.9\times$ faster, at a fraction of the quality loss of plain DDIM at the same
step count.

\subsection{Ablation 1 (causal): distillation objective vs.\ naive eps-loss}
To prove the gain stems from the distillation \emph{objective} and not from the added
LoRA capacity, we train an identical LoRA (same rank, same trainable parameters) with
the standard noise-prediction loss (Eq.~\eqref{eq:ddpm-loss}) at low steps.

\begin{table}[H]
\centering
\caption{Causal control: naive eps-loss LoRA vs.\ distillation, matched $K$ and rank.}
\label{tab:causal}
\begin{tabular}{crr}
\toprule
$K$ & Naive-eps LoRA $\Delta$FID & Distillation $\Delta$FID \\
\midrule
8 & $+0.96$ & \textbf{+8.9} \\
4 & $-3.51$ (worse) & \textbf{+34.9} \\
\bottomrule
\end{tabular}
\end{table}

The naive control yields $\approx\!0$ improvement at $K=8$ and is \emph{worse} than
plain DDIM at $K=4$, while distillation yields large gains. This isolates the
trajectory-matching objective as the direct cause of the improvement.

\subsection{Ablation 2: LoRA rank ($K=8$)}
FID: $r{=}2$: $38.32$, $r{=}4$: $38.48$, $r{=}8$: $37.67$. Rank barely matters --- even
rank~2 captures almost all of the benefit, reinforcing the parameter-efficient claim.

\subsection{Ablation 3: distillation state source ($K=8$)}
Forward-diffused states only: $38.48$; mixing teacher-trajectory states ($p=0.5$):
$38.24$ --- essentially tied. Conditioning on generated trajectory states does
\emph{not} help; forward-diffusion states are sufficient. We report this as a negative
result and do not rely on trajectory-state mixing in the final method.

% ==================================================================
\section{Limitations}
% ==================================================================
\begin{itemize}
  \item \textbf{Weak absolute baseline.} Under the single-GPU budget the baseline
        reaches FID $\approx 29$; samples show global structure but limited detail.
        Because all comparisons are relative to the same teacher, the conclusions about
        the novelty hold regardless of absolute quality.
  \item \textbf{Breakdown at $K=1$.} The method does not recover single-step sampling on
        this baseline.
  \item \textbf{FID sample-size sensitivity.} We evaluate all headline comparisons at a
        fixed \num{10000} images. The $K=4$ result ($+34.9$) is well beyond this noise;
        the $K=8$ result ($+8.9$) is smaller and would benefit from multiple seeds.
  \item \textbf{Linear $\beta$-schedule.} A cosine schedule could improve baseline
        quality and is left as future work.
\end{itemize}

% ==================================================================
\section{Conclusion}
% ==================================================================
We presented a parameter-efficient few-step distillation of a DDPM using LoRA adapters
and a trajectory-matching objective. On CIFAR-10 the method improves few-step generation
quality at a matched step budget (up to $+34.9$ FID at 4 steps), partially restores the
diversity lost by aggressive step reduction, and adds no inference overhead, yielding
$\sim\!10$--$19\times$ speed-ups over the DDIM-100 baseline (over $100\times$ vs.\ full
DDPM). A naive-eps control confirms that the distillation objective --- not the extra
parameters --- is responsible for the gain. Future work includes a cosine schedule,
min-SNR-weighted distillation, and rank scaling for the extreme $K=2/1$ regime.

% ==================================================================
\begin{thebibliography}{9}
\bibitem{ho2020ddpm} J.~Ho, A.~Jain, P.~Abbeel.
  \emph{Denoising Diffusion Probabilistic Models.} NeurIPS 2020.
\bibitem{song2021ddim} J.~Song, C.~Meng, S.~Ermon.
  \emph{Denoising Diffusion Implicit Models.} ICLR 2021.
\bibitem{salimans2022progressive} T.~Salimans, J.~Ho.
  \emph{Progressive Distillation for Fast Sampling of Diffusion Models.} ICLR 2022.
\bibitem{hu2021lora} E.~Hu, Y.~Shen, P.~Wallis, et al.
  \emph{LoRA: Low-Rank Adaptation of Large Language Models.} 2021.
\bibitem{kynkaanniemi2019precision} T.~Kynk\"a\"anniemi, T.~Karras, S.~Laine,
  J.~Lehtinen, T.~Aila. \emph{Improved Precision and Recall Metric for Assessing
  Generative Models.} NeurIPS 2019.
\end{thebibliography}

\end{document}
